\documentclass[crop,tikz,11pt]{standalone}
\usepackage{geometricfigures}
\usepackage{amsmath,amssymb}
\usepackage{pgfplots}
\pgfplotsset{compat=1.18}
\usetikzlibrary{arrows.meta,decorations.markings,patterns,intersections}
\usepgfplotslibrary{fillbetween}

\definecolor{lib}{HTML}{3B75AF}
\definecolor{rot}{HTML}{C1701E}
\definecolor{sep}{HTML}{B03030}
\definecolor{net}{HTML}{4E9A4E}

% M is drawn through the projection R^4 -> R^3, x |-> (x_1,x_2,x_2x_3-x_1x_4),
% which restricts to M as (sin t, cos t, p_theta): the round cylinder. That is
% put on the page with u = theta - pi as
%   X(u) = R sin u,   Y(u,p) = -c R cos u + s p,
% so u=0 (the stable equilibrium) is at the front and u=+-pi at the back.
\newcommand{\cylR}{2.25}   % radius on the page
\newcommand{\cylC}{0.40}   % foreshortening of the circle
\newcommand{\cylS}{0.80}   % scale of the p_theta axis
\newcommand{\cylP}{3.25}   % where the picture cuts the (unbounded) cylinder
\pgfmathdeclarefunction{cylx}{1}{\pgfmathparse{\cylR*sin(#1)}}
\pgfmathdeclarefunction{cyly}{2}{\pgfmathparse{-\cylC*\cylR*cos(#1)+\cylS*(#2)}}
\pgfmathdeclarefunction{cylorb}{2}{\pgfmathparse{sqrt(max(0,2*((#2)+cos(#1))))}}

\begin{document}
\begin{tikzpicture}[>=Stealth,
  hidden/.style={dash pattern=on 2.2pt off 2.2pt,opacity=0.38}]

  % --- the cylinder: body, the two cut rims, the silhouette -----------------
  \fill[fg!4!bg]
    plot[domain=-90:90,samples=60,smooth] ({cylx(\x)},{cyly(\x,-\cylP)})
    -- plot[domain=90:270,samples=60,smooth] ({cylx(\x)},{cyly(\x,\cylP)})
    -- cycle;
  \draw[fg!25!bg,densely dotted]
    plot[domain=-180:180,samples=140,smooth] ({cylx(\x)},{cyly(\x,\cylP)});
  \draw[fg!25!bg,densely dotted]
    plot[domain=-180:180,samples=140,smooth] ({cylx(\x)},{cyly(\x,-\cylP)});
  \draw[fg!25!bg] ({-\cylR},{-\cylS*\cylP}) -- ({-\cylR},{\cylS*\cylP});
  \draw[fg!25!bg] ({\cylR},{-\cylS*\cylP}) -- ({\cylR},{\cylS*\cylP});

  % --- far side of the cylinder, dashed ------------------------------------
  \draw[rot,thick,hidden] plot[domain=90:270,samples=110,smooth] ({cylx(\x)},{cyly(\x,cylorb(\x,1.55))});
  \draw[rot,thick,hidden] plot[domain=90:270,samples=110,smooth] ({cylx(\x)},{cyly(\x,-cylorb(\x,1.55))});
  \draw[rot,thick,hidden] plot[domain=90:270,samples=110,smooth] ({cylx(\x)},{cyly(\x,cylorb(\x,2.60))});
  \draw[rot,thick,hidden] plot[domain=90:270,samples=110,smooth] ({cylx(\x)},{cyly(\x,-cylorb(\x,2.60))});
  \draw[sep,very thick,hidden] plot[domain=90:270,samples=140,smooth] ({cylx(\x)},{cyly(\x,cylorb(\x,1.0))});
  \draw[sep,very thick,hidden] plot[domain=90:270,samples=140,smooth] ({cylx(\x)},{cyly(\x,-cylorb(\x,1.0))});
  % the outermost librating loop reaches a little past the silhouette
  \draw[lib,thick,hidden]
    plot[domain=90:116.74,samples=34,smooth] ({cylx(\x)},{cyly(\x,cylorb(\x,0.45))})
    -- plot[domain=116.74:90,samples=34,smooth] ({cylx(\x)},{cyly(\x,-cylorb(\x,0.45))});
  \draw[lib,thick,hidden]
    plot[domain=-90:-116.74,samples=34,smooth] ({cylx(\x)},{cyly(\x,cylorb(\x,0.45))})
    -- plot[domain=-116.74:-90,samples=34,smooth] ({cylx(\x)},{cyly(\x,-cylorb(\x,0.45))});

  % --- near side ------------------------------------------------------------
  % librating: three level sets; the innermost is shaded with the disk it bounds
  \fill[lib!14!bg]
    plot[domain=-53.13:53.13,samples=60,smooth] ({cylx(\x)},{cyly(\x,cylorb(\x,-0.60))})
    -- plot[domain=53.13:-53.13,samples=60,smooth] ({cylx(\x)},{cyly(\x,-cylorb(\x,-0.60))}) -- cycle;
  \draw[lib,thick]
    plot[domain=-53.13:53.13,samples=60,smooth] ({cylx(\x)},{cyly(\x,cylorb(\x,-0.60))})
    -- plot[domain=53.13:-53.13,samples=60,smooth] ({cylx(\x)},{cyly(\x,-cylorb(\x,-0.60))}) -- cycle;
  \draw[lib,thick]
    plot[domain=-75.52:75.52,samples=80,smooth] ({cylx(\x)},{cyly(\x,cylorb(\x,-0.25))})
    -- plot[domain=75.52:-75.52,samples=80,smooth] ({cylx(\x)},{cyly(\x,-cylorb(\x,-0.25))}) -- cycle;
  \draw[lib,thick] plot[domain=-90:90,samples=80,smooth] ({cylx(\x)},{cyly(\x,cylorb(\x,0.45))});
  \draw[lib,thick] plot[domain=-90:90,samples=80,smooth] ({cylx(\x)},{cyly(\x,-cylorb(\x,0.45))});
  % separatrix
  \draw[sep,very thick] plot[domain=-90:90,samples=110,smooth] ({cylx(\x)},{cyly(\x,cylorb(\x,1.0))});
  \draw[sep,very thick] plot[domain=-90:90,samples=110,smooth] ({cylx(\x)},{cyly(\x,-cylorb(\x,1.0))});
  % rotating
  \draw[rot,thick] plot[domain=-90:90,samples=80,smooth] ({cylx(\x)},{cyly(\x,cylorb(\x,1.55))});
  \draw[rot,thick] plot[domain=-90:90,samples=80,smooth] ({cylx(\x)},{cyly(\x,-cylorb(\x,1.55))});
  \draw[rot,thick] plot[domain=-90:90,samples=80,smooth] ({cylx(\x)},{cyly(\x,cylorb(\x,2.60))});
  \draw[rot,thick] plot[domain=-90:90,samples=80,smooth] ({cylx(\x)},{cyly(\x,-cylorb(\x,2.60))});

  % --- the two equilibria ---------------------------------------------------
  \fill[sep] (0,{\cylC*\cylR}) circle (1.7pt);
  \node[sep,font=\footnotesize,anchor=east,inner sep=1.5pt,
        fill=bg,fill opacity=0.82,text opacity=1] at (-0.13,{\cylC*\cylR}) {$\theta{=}0$};
  \fill[lib] (0,{-\cylC*\cylR}) circle (1.7pt);
  \node[lib,font=\footnotesize,anchor=west,inner sep=1.5pt,
        fill=bg,fill opacity=0.82,text opacity=1] at (0.13,{-\cylC*\cylR}) {$\theta{=}\pi$};

  % name the two separatrix branches where they cross the left silhouette
  \node[sep,font=\footnotesize,anchor=east,inner sep=1.5pt,
        fill=bg,fill opacity=0.82,text opacity=1]
    at (-2.30,{\cylS*1.4142}) {$\ell_+$};
  \node[sep,font=\footnotesize,anchor=east,inner sep=1.5pt,
        fill=bg,fill opacity=0.82,text opacity=1]
    at (-2.30,{-\cylS*1.4142}) {$\ell_-$};

  % --- the two directions ---------------------------------------------------
  \draw[->,fg!50!bg] (2.85,0.15) -- (2.85,1.35)
    node[above,font=\small,fg,inner sep=1pt] {$p_\theta$};
  \draw[->,fg!50!bg] plot[domain=-34:34,samples=26,smooth] ({cylx(\x)},{cyly(\x,-\cylP)});
  \node[font=\small,fg!70!bg,anchor=north] at (0,{-\cylC*\cylR-\cylS*\cylP-0.05}) {$\theta$};

  % --- what the picture says ------------------------------------------------
  \node[anchor=east,text width=4.5cm,align=right,font=\small] at (-3.3,2.3)
    {$\mathcal{M}\subset\mathbb{R}^4$: the solution manifold, the image of the cotangent lift.
     It is a cylinder $S^1\times\mathbb{R}$, drawn here in $\mathbb{R}^3$ --- see the caption.};
  \node[anchor=east,text width=4.5cm,align=right,font=\small] at (-3.3,-1.35)
    {\textcolor{lib}{\textbf{librating}}, $H<1$: a loop that \emph{bounds} in $\mathcal{M}$.
     $\oint p_\theta\,d\theta$ is the symplectic area of the disk it encircles, so every symplectic chart
     must report the same number.};
  \node[anchor=west,text width=4.5cm,align=left,font=\small] at (3.3,2.15)
    {\textcolor{rot}{\textbf{rotating}}, $H>1$: a loop that runs \emph{around} the cylinder.
     It bounds no surface in $\mathcal{M}$; its action uses the canonical one-form
     $p_\theta\,d\theta$, independently of the chart.};
  \node[anchor=west,text width=4.5cm,align=left,font=\small] at (3.3,-2.05)
    {\textcolor{sep}{\textbf{separatrix}}, $H=1$: the two branches $\ell_\pm$, labelled by the
     sign of $p_\theta$, each with action $+8$ along the flow. They meet at the unstable
     equilibrium $\theta=0$.};
\end{tikzpicture}
\end{document}
